\subsubsection{Dépendance structurelle de la coordonnée électronique}

L'équation électronique exprime une composition hiérarchique. Son premier
facteur, \(\sqrt3/4\), appartient à la structure matricielle centrale et fixe
une normalisation géométrique. L'exposant \(\varphi/\pi^2\) utilise les
coordonnées régionales déjà engendrées dans l'orientation indiquée de l'opérateur
de lecture surfacique--volumique. Le terme \(22\alpha^3\) incorpore le couplage
dodécaphasique à la correction cubique. Le facteur \(1+15\Delta_4\) conserve
la normalisation torsionnelle globale. Enfin, le quotient des sections
d'action agit avec l'exposant déterminé par les opérateurs de lecture du retour.
Chaque opération possède ainsi un antécédent distinct au sein de la même
construction.

L'articulation devient plus précise lorsqu'on considère ses dépendances.
La norme centrale peut être calculée avant l'évaluation des coordonnées
régionales. L'exponentielle et la correction cubique requièrent ces
coordonnées et la clôture de \(\alpha\). La mémoire multiplicative requiert
en outre les sections d'action et la trajectoire de \(108\) positions.
Connaître l'expression basale n'équivaut donc pas à avoir évalué la
coordonnée complète. La succession d'opérations n'ajoute pas une masse
externe au point \((5,5)\) ; elle développe les opérateurs de lecture attribués à sa
section persistante.

La distinction entre échelle et structure fournit également une
interprétation de la comparaison physique. La masse est la réalisation
dimensionnelle de la coordonnée complète, tandis que le centre, l'involution
et la lacune décrivent les relations qui la produisent. Un changement commun
d'unité modifie les coordonnées dimensionnelles, mais laisse inchangés
les rapports sans dimension et les identités de transport démontrées.
C'est en ce sens que l'électron conserve son caractère structurel
lorsqu'il s'exprime par rapport à une échelle chronologique.
